\documentclass[10pt,a4paper]{amsart}
\usepackage[margin=19mm]{geometry}
\usepackage{amsmath,amssymb,mathtools}
\usepackage[scr=rsfs,cal=euler]{mathalfa}
\usepackage{xfrac}
\usepackage[unicode,hidelinks,bookmarks=false]{hyperref}
\newcommand{\ie}{i.e.\ }
\newcommand{\cf}{cf.\ }
\newcommand{\resp}{respectively\ }
\newcommand{\Subsection}{\subsection}
\providecommand{\nobreakdash}{\nobreak\hbox{-}}
\setlength{\emergencystretch}{2em}
\multlinegap0pt
\usepackage{amssymb}
\usepackage{tikz}
\usepackage{xcolor}
\usepackage{float}
\usepackage[shortlabels]{enumitem}
\newtheorem{theorem}{Theorem}
\newtheorem{proposition}{Proposition}
\DeclareMathOperator{\Aut}{Aut}
\DeclareMathOperator{\id}{id}
\title{On Chamber-regular $\tilde C_2$-Lattices}
\author{Franziska Stamer, Thomas Titz Mite}
\date{Source: arXiv:2511.08312v2 (CC BY 4.0)}
\begin{document}
\maketitle

\noindent\emph{Source and licence.} On Chamber-regular $\tilde C_2$-Lattices. Version arXiv:2511.08312v2 (CC BY 4.0), distributed by arXiv under the Creative Commons Attribution 4.0 International licence, http://creativecommons.org/licenses/by/4.0/, as recorded in the arXiv licence metadata for this exact version and shown on the abstract page https://arxiv.org/abs/2511.08312v2. Full licence notice: \texttt{licenses/arxiv-2511.08312-CC-BY-4.0.txt}.\par\smallskip

\noindent\textit{Editorial scope.} Counted unit: the article's proposition prop:action\_rigidity (source0.tex lines 535--538). The article's complete original proof of it is delivered with it (source0.tex lines 540--580, one \texttt{proof} environment of 41 lines). No mathematics is written, repaired or re-derived: the statement and the proof are the article's own lines, in the article's own notation, and no retained line is substituted or re-flowed.  Retained uncounted as the source-local context the proof needs: lines 528--531.  Nothing was omitted from any retained span, so the package carries no omission record.  No retained line contains a citation, so the wrapper carries no bibliography and no bibliography entry.  No retained line contains a figure environment, an image-inclusion command or drawing code, so no image is delivered and the package declares no asset dependency.  Editorial formatting only, in addition: an AMS \texttt{amsart} preamble that restates the article's own declarations and macros used by the retained lines, namely the environments \texttt{theorem}, \texttt{proposition} on separate counters, together with the standard packages those lines need.  The retained environments print in this document as Theorem 1, Proposition 1 in order of appearance, whereas the article prints the same items with the numbering of its own full document.  The complete native source of the article as distributed in the arXiv e-print for this exact version is bundled unchanged as \texttt{sources/188880/source0.tex}.
    \begin{theorem}\label{thm:qi_rigidity}
        Let $X_1$ and $X_2$ be locally finite Euclidean buildings of dimension at least two. Let $f : X_1 \to X_2$ be a quasi-isometry. Then there
        exists unique isometry $\bar f:X_1 \to X_2$ that is at bounded distance from $f$.
    \end{theorem}

    \begin{proposition}\label{prop:action_rigidity}
        Let $\Gamma$ be an abstract group and let $X_1,X_2$ be localy finite Euclidean buildings of dimension at least two.
        Let $\phi_1 : \Gamma \to \Aut(X_1)$ and $\phi_2 : \Gamma \to \Aut(X_2)$ be injective maps, such that the images are uniform lattices on $X_1$ and $X_2$ respectively. Then the actions of $\phi_1(\Gamma)$ and $\phi_2(\Gamma)$ are isomorphic, i.e. there exists an isomorphism $\theta: X_1\to X_2$ such that $\theta \circ \phi_1(g) = \phi_2(g)\circ \theta$ for all $g \in \Gamma$.
    \end{proposition}

    \begin{proof}
        By the \v{S}varc--Milnor lemma, $\Gamma$ is quasi-isometric to both $X_1$ and $X_2$, so $X_1$ and $X_2$ are isomorphic by Theorem \ref{thm:qi_rigidity}. 
        So without loss of generality, we can assume that $X_1=X_2$ and denote the building just by $X$.
        Choose a base point $o \in X$. Define $\tilde \phi_i : \Gamma \to X, g \mapsto \phi_i(g)(o)$ for $i \in \{1,2\}$.
        We slightly perturb the image of $\tilde \phi_2$, such that it becomes injective and $\tilde \phi_2(g)$ stays close to $\phi_2(g)(o)$.
        Let $\kappa_2 : X\to \Gamma$ be a quasi-inverse such that $\kappa_2 \circ \tilde \phi_2 = \id_\Gamma$.
        Then $\tilde \phi_1 \circ \kappa_2 : X \to X$ is a quasi-isometry. By Theorem \ref{thm:qi_rigidity}, there is a unique $\theta\in\Aut(X)$ at finite distance $D$:
        \[
            d(\tilde \phi_1 \circ \kappa_2(x), \theta(x)) < D \quad \text{for all } x\in X.
        \]
        We claim that for all $g \in \Gamma$, we have that $\theta \circ \phi_2(g) = \phi_1(g) \circ\theta$. In particular the lattices $\phi_1(\Gamma)$ and $\phi_2(\Gamma)$ are conjugated in $\Aut(X)$.
        We show that $\theta \circ \phi_2(g)$ and $\phi_1(g) \circ \theta$ are at finite distance and therefore equal by Theorem \ref{thm:qi_rigidity}.
        Since $\phi_2(\Gamma).o$ is cobounded, it suffices to show that for all $h \in \Gamma$ the distances
        \[
            d((\theta \circ \phi_2(g)(\phi_2(h)(o)), (\phi_1(g)\circ \theta)(\phi_2(h)(o)))
        \]
        are uniformly bounded.
        We have 
        \[
            (\theta \circ \phi_2(g))(\phi_2(h)(o)) = \theta \circ \phi_2(gh)(o),
        \]
        and this point is at distance at most $D$ from the following point
        \[
            \tilde \phi_1 \circ \kappa_2 \circ \phi_2(gh)(o)
            = \tilde \phi_1\circ \kappa_2 \circ \tilde \phi_2(gh) 
            = \tilde \phi_1 (gh)
            = \phi_1(gh)(o).
        \]
        On the other hand, the point
        \[
            (\phi_1(g) \circ \theta)(\phi_2(h)(o))
        \]
        is at distance at most $D$ from 
        \begin{align*}
            \phi_1(g) \circ \tilde \phi_1 \circ \kappa_2 \circ \phi_2(h) (o)
            &= \phi_1(g) \circ \tilde \phi_1 \circ \kappa_2 \circ \tilde \phi_2(h)
            &&= \phi_1(g) \circ \tilde \phi_1(h) \\
            &= \phi_1(g) \circ \phi_1(h)(o)
            &&\hspace{0.2mm}= \phi_1(gh)(o).
        \end{align*}
    \end{proof}

\end{document}
